% Auto-generated from Chapter-03-The-Robbery.md
% Source: C:\Users\PCGAME\Desktop\story&universes\chain-weapon-story\finished-manuscript\manuscriptbaseforchapters\Chapter-03-The-Robbery.md

\chapter{The Robbery \& First Kill}
\renewcommand{\CurrentChapterNum}{3}
\POV{\glyphAisen}{Aisen}

The night it happened, Aisen was helping his father count coins.

This was not unusual. On most nights when business had closed, his father brought the register to the private room behind the bar and Aisen, now eight years old, was allowed to sit quietly and watch—not yet counting, but learning the rhythm of it. Coins sorted by weight, by mint mark, by trustworthiness. Some were worn nearly smooth from passage between hands. Some bore the marks of forgeries that careful persons had attempted and that his father rejected with the expertise of someone who'd handled thousands of coins and understood their density, their sound, their feel.

His father was teaching him this without quite teaching him. Simply by allowing him presence, by demonstrating precision, by handling each coin as if it mattered—which, of course, it did.

The breach was not quiet.

It came as a sound first—the bar door, the one that faced the street, being forced rather than opened. The hinges didn't give way; whoever pushed was desperate, not strong. But desperation could accomplish what strength sometimes couldn't. The door swung inward, and the night wind came with it, bringing the smell of cold and wet earth.

His father's hand stilled on the coin he'd been examining.

"Stay here," his father said quietly to Aisen. Not a question, not a request. A statement of what would happen. Aisen had learned to recognize this tone, had learned that when his father used it, obedience was the only sensible response.

Aisen pressed himself against the wall of the small counting room and listened.

He heard voices—his father's calm and steady, asking what was wanted. And someone else's voice, younger maybe, definitely panicked. "Coins!" the voice said, then something that made no sense: "Don't make this hard!"

Don't make it hard? The person breaking into a tavern wanted it to be easy? That seemed upside down to Aisen. Everything about this was backwards.

He heard his father say, "There's not much here. You're welcome to what's in the main register, but if you want the night's full take, it's in the back room which is where I was just headed."

Aisen's breath caught. His father, who never lied about small things, was lying in a way that was clearly calculated. The main register held perhaps a third of the night's coins. The back room—this room, where Aisen was hiding—held the rest. So his father was offering the robber a choice: take what was easy, or come further in and take what was harder.

Why would he do that?

Then Aisen understood, the way children understand things. His father was moving the threat. Moving it toward territory where he had control. This was the same skill he used in the tavern every day—knowing position, knowing advantage, knowing the geography of influence.

Aisen heard footsteps, more than one. The robber wasn't alone then. Or possibly his father was wrong and it was something else entirely.

The voice came again, panicked and therefore dangerous: "Where's the rest?"

"Back room," his father said simply. "But understand, you're committed now. Robbery of a closed tavern at night—that's a different crime from theft of loose coins. You want to keep thinking this through, or you want the coins and to move on?"

Aisen's mind was working quickly now, the way his observations had trained it to work. The robber was young, was panicked, was desperate but not experienced. His father was trying to scare him into taking the easy robbery and leaving. But something about this was wrong. Something about the situation had a quality that made his father alert.

He heard the robber say something low, something that might have been agreement or might have been refusal—Aisen couldn't quite catch it.

Then his father's voice again, "The door to the back room is down this hall. But I should tell you—there's someone back there. I'm not going to ask him to come out. I'm just telling you so you have information."

Aisen's heart was beating very fast now. He was about to be discovered. He was about to be seen in the middle of what was already a crime. His father had just revealed his presence without protecting him, and that meant—what did that mean?

The sounds were getting closer.

Aisen, without thinking about it, without planning it, moved. He grabbed the brass scale that sat on the counting table—heavy, substantial, clearly useful. It was the first thing his hands found, and his young mind did a calculation: if the robber came through the door, the element of surprise was gone, but the element of assault was there.

The robber appeared in the doorway.

He was young. That was the first thing Aisen registered. Not much older than Harald, who was maybe twenty-five. This man was maybe sixteen or seventeen, thin in a way that suggested not poverty exactly but lack—lack of food,  lack of sleep, lack of certainty about his place in the world. He held a knife that looked both functional and desperate, the kind of knife made for farming work rather than killing.

And he was off-balance.

Later, Harald would tell him that the moment of decision happens in milliseconds, that the people who survive combat are the people who understand their environment in ways that others don't. But in this moment, Aisen didn't think about survival. He simply understood that the robber, coming through the doorway into a room that wasn't quite what he expected, with a child present and a man in the hallway still blocking escape, was experiencing a moment of disorientation.

Aisen threw the scale.

It didn't connect with violence. It connected with surprise. The scale hit the robber's shoulder and clattered to the ground, and in that instant of distraction, everything else accelerated.

His father was there suddenly, moving with economical purpose. Aisen would later learn that this was Harald's influence, that his father had learned to move this way in preparation for chaos, but in the moment all Aisen understood was that his father was stopping the robbery the way he stopped difficult customers—through positioning and intent rather than force.

But then there was another sound. Running feet. Someone else, entering from the front of the tavern. And Aisen realized with a clarity that was almost mathematical: the robbery was a distraction. While the desperate young man drew attention, someone else was--what? Stealing what?

Then the second figure appeared, and Aisen understood he was wrong again.

It was Harald.

He was in his guard uniform, which meant he'd arrived to begin his evening lodging and had found the door breached. And something in him, something that had been trained into him by whatever formed young guards in border towns, recognized what was happening and responded.

The young robber, seeing Harald, made a choice.

He lunged at Aisen's father with the knife.

And Harald moved.

Later, people would ask Aisen to describe this moment, and he would struggle because what happened was efficient to the point of being almost disappointing. There was no grand combat. Harald didn't engage the robber so much as redirect him, the way a river redirects a stone. The knife that had been aimed at Aisen's father found empty air as the robber, confused by his own momentum and by Harald's presence, stumbled and fell.

But falling was not the end of motion. The robber, desperate and now truly panicked, turned the knife on himself—not consciously, Aisen thought, but in the desperation of being trapped. He saw the opportunity that wasn't there, saw the path that didn't exist, and his own momentum carried him onto his own blade.

It was not quick. It was not painless. It was the kind of death that made Aisen's eight-year-old mind recoil from what it was seeing because some things were too real to be processed in the moment.

Harald caught the falling man and eased him to the ground with the care of someone who'd seen this before and understood that even a failed robber deserved dignity in their ending.

It took perhaps three seconds. It felt like hours.

His father moved immediately to Aisen.

"Are you hurt?" Not "Are you afraid?" or "Are you shocked?" but the practical question, the one that mattered first.

Aisen shook his head, though his hands were shaking and his vision was doing something strange and narrow.

His father held him—not speaking, just holding—while Harald worked in the hallway to secure the tavern and ensure no one else was using the robbery as cover for greater theft.

When Harald came back to the counting room, he was moving carefully, aware of the child present but not treating him like a child who couldn't understand. He knelt beside the dead robber and checked for pulse and breath and confirmed what was already obvious: the man was gone.

"It resolved in our favor," Harald said quietly to Aisen's father. "The robber died. You understand."

Aisen's father nodded slightly.

Harald looked at Aisen directly. "You threw the scale," he said. It wasn't a question.

Aisen managed to nod.

"You noticed he was off-balance," Harald continued. "You understood that a moment of confusion might be valuable. You acted on that understanding." Harald paused. "You thought. Most people freeze. You thought."

Something in Aisen's chest, something that had been compressed since the breach, began to release slightly.

"His father will have to arrange burial," Aisen's father said flatly. He was already thinking ahead, was already moving into the practical matters that followed violence. "I'll need to report this to the guard. This is a self-defense case, but it needs documentation."

"I'll handle the documentation," Harald said. His voice carried the weight of someone who had the authority to make such a statement. "He was attempting armed robbery. You defended. The boy threw an object that created a moment of opportunity. The robber fell on his own blade during the struggle. It's simple."

Over the following hours, while the city guard came and documented and eventually removed the body, Aisen realized two things.

First, that his observation about the robber being desperate rather than confident had been correct. The young man was already known to the guard—a failed apprentice, a person who'd lost his way, a community failure that had finally ended. In another world, he might have taken his own life more deliberately. In this world, his desperation had given him to his own violence.

Second, that his father had known this was possible. His father lived at the intersection of roads where desperate people sometimes came. His father, through observation and information, had understood that robbery might come. His father had been preparing without apparent preparation, the way his father did everything.

When the guard had left and the counting room had been cleaned and the body had been removed, Harald sat across from Aisen and his father.

"I'd like to teach him properly," Harald said. "What he did today was instinct. But instinct can be refined. Can be made reliable. If you'll permit it."

Aisen's father looked at Aisen. Not asking permission but gauging something—whether Aisen was too damaged by what he'd witnessed, whether this was the right moment, whether Aisen had something in him that wanted this.

Aisen met his father's eyes and said nothing. But something in his face must have conveyed assent because his father turned to Harald and nodded.

"He's yours," his father said. "But carefully. He's still young."

"I know," Harald said. "That's why it matters. If we teach him to think early, he'll continue thinking when others panic."

That night, Aisen lay in his small bed above the kitchen and understood that something fundamental had changed. He had seen what death looked like. He understood that his body could be used to make choices. He understood that the knife falling and the blood and the finality of it were all part of being alive in a place where violence sometimes came through the door and needed to be answered.

He was eight years old. He would carry this night differently than other children carried their eight-year-old selves. But he would carry it with someone who understood—with Harald, who would teach him that thinking was the most valuable thing he possessed, that observation was a form of power, that the moment between stimulus and response was where all real choice happened.

He was alive. The robber was dead. His father was safe. And somewhere in the compass of that reality, Aisen had found both his fear and his purpose, had found both the fragility of life and its possibility, had found the beginning of the person he would become.

By morning, the tavern had been cleaned and reset. The business of \emph{The Crossing} continued. Customers came and went. His father served them and took their coins and demonstrated the ordinary steady that made safety possible.

And Aisen, watching from his habitual places, began the work of learning to make choices.